% 原书第 5—7 页；第 1.3 节、图 1.1—1.2、原理 1.1。
\FloatBarrier
\section{因果建模与学习}
\label{sec:causal-modeling-learning}

\textbf{因果建模（causal modeling）}从另一种可以说更为根本的结构出发。一个因果结构能够导出一个概率模型，但它还包含后者所没有的额外信息，参见第 1.4 节的例子。按照本书的术语约定，\textbf{因果推理（causal reasoning）}是指从因果模型中推出结论的过程，这类似于概率论使我们能够对随机试验的结果进行推理。然而，由于因果模型包含的信息多于概率模型，因果推理也就比概率推理更有力：它使我们能够分析干预或分布变化所产生的影响。

正如统计学习是概率论的逆问题，我们也可以考虑，如何根据因果结构在经验层面上的表现来推断因果结构本身。这里的经验信息可以完全来自观测，也可以包括干预下的数据，例如随机化试验数据，或者分布变化下的数据。研究者用不同术语来指称这类问题，包括\textbf{结构学习（structure learning）}和\textbf{因果发现（causal discovery）}。另一个密切相关的问题是：原则上，因果结构的哪些部分能够从联合分布中推断出来？我们将其称为\textbf{结构可识别性（structure identifiability）}问题。与第 1.2 节介绍的标准统计学习问题不同，即使完全知道分布 $P$，这个问题也不会变得容易，我们仍需作出额外假设，参见第 2、4、7 章。

\begin{figure}[htbp]
\centering
\begin{tikzpicture}[causal,x=1cm,y=1cm,font=\small,
 modelbox/.style={draw,rounded corners=2mm,align=center,minimum width=2.6cm,minimum height=.8cm,inner sep=3mm}]
 \node[modelbox] (f11causal) at (1.4,3.1) {因果模型};
 \node[modelbox,minimum width=3.5cm,minimum height=1.4cm] (f11causalobs) at (11.2,3.1)
   {观测与结果，\\包括变化与干预};
 \node[modelbox] (f11prob) at (1.4,0) {概率模型};
 \node[modelbox,minimum width=3.5cm,minimum height=1.0cm] (f11probobs) at (11.2,0)
   {观测与结果};
 \draw[ed] ([yshift=2mm]f11causalobs.west)
   to[bend right=11] node[above=2mm] {因果学习}
   ([yshift=2mm]f11causal.east);
 \draw[ed] ([yshift=-2mm]f11causal.east)
   to[bend right=11] node[below=2mm] {因果推理}
   ([yshift=-2mm]f11causalobs.west);
 \draw[ed] ([yshift=2mm]f11probobs.west)
   to[bend right=11] node[above=2mm] {统计学习}
   ([yshift=2mm]f11prob.east);
 \draw[ed] ([yshift=-2mm]f11prob.east)
   to[bend right=11] node[below=2mm] {概率推理}
   ([yshift=-2mm]f11probobs.west);
 \draw[->,dashed,gray] (f11causal.south) -- node[left=2mm,text=black] {包含} (f11prob.north);
 \draw[->,dashed,gray] (f11causalobs.south) -- node[right=2mm,text=black] {包含} (f11probobs.north);
\end{tikzpicture}
\caption{本书描述各类概率推断问题（下方）与因果推断问题（上方）时所使用的术语，见第 1.3 节。注意，本书所说的“推断”同时包括学习和推理。}
\label{fig:1.1}
\end{figure}


不过，这一困难不应使我们忽略通常统计问题中的不适定性依然存在。因此，在因果问题中，同样有必要关注函数类的容量，例如使用加性噪声模型，见第 4.1.4 节。只不过，我们现在试图估计的是比概率结构更加丰富的结构，因而又多了一层困难。我们将这一整体问题称为\textbf{因果学习（causal learning）}。图 1.1 总结了上述问题和模型之间的关系。

要从观测分布中学习因果结构，我们需要理解因果模型与统计模型之间的联系。第 4 章和第 7 章将重新讨论这一问题，这里先给出一个例子。人们熟知的一个说法是：“相关并不意味着因果关系。”换言之，仅靠统计性质不能确定因果结构。另一个较少为人所知的观点则是：虽然我们不能推断出具体的因果结构，但至少可以从统计相依关系推断出因果联系的存在。Reichenbach~[1956] 最先认识到了这一点。下面将他的这一认识表述为一个原理，另见图 1.2。\footnote{为使表述清晰，我们将某些重要假设称为原理。本书并不始终将这些假设视为当然成立；在这个意义上，它们不是公理。}

\begin{figure}[htbp]
\centering
\begin{tikzpicture}[causal,x=1cm,y=1cm]
 \node[vtx] (f12aX) at (0,0) {$X$};
 \node[vtx] (f12aY) at (2.6,0) {$Y$};
 \node[latent] (f12aZ) at (1.3,1.1) {$Z$};
 \draw[edd] (f12aZ)--(f12aX);
 \draw[edd] (f12aZ)--(f12aY);
 \node[vtx] (f12bX) at (5.0,0) {$X$};
 \node[vtx] (f12bY) at (7.0,0) {$Y$};
 \draw[ed] (f12bX)--(f12bY);
 \node[vtx] (f12cX) at (9.4,0) {$X$};
 \node[vtx] (f12cY) at (11.4,0) {$Y$};
 \draw[ed] (f12cY)--(f12cX);
\end{tikzpicture}
\caption{Reichenbach 的共同原因原理在统计性质与因果结构之间建立了联系。两个可观测变量 $X$ 和 $Y$ 之间的统计相依关系，表明它们受到某个变量 $Z$ 的共同影响；$Z$ 通常称为混杂因子（左图）。这里，$Z$ 也可以与 $X$ 或 $Y$ 中的一个重合，此时图形简化为中图或右图。该原理还指出，给定 $Z$ 时，$X$ 与 $Y$ 统计独立。图中的箭头表示直接因果关系，参见第 3 章和第 6 章。}
\label{fig:1.2}
\end{figure}


\begin{principle}[Reichenbach 共同原因原理]
\label{principle:1.1}
如果两个随机变量 $X$ 与 $Y$ 统计相依（$X\nindep Y$），那么就存在第三个变量 $Z$，对两者都产生因果影响。作为特殊情形，$Z$ 可以与 $X$ 或 $Y$ 中的一个重合。此外，变量 $Z$ 能够将 $X$ 与 $Y$ 彼此\textbf{屏蔽（screening off）}：给定 $Z$ 后，两者独立，即 $X\indep Y\mid Z$。
\end{principle}

在实际问题中，相依关系也可能由于共同原因原理未提及的其他原因而出现，例如：
\begin{enumerate}[label=(\arabic*)]
 \item 我们观测到的随机变量，实际上已经以其他变量为条件；这种条件化往往由选择偏倚隐含地引入。我们将在后文重新讨论这一问题，参见注 6.29。
 \item 随机变量只是\emph{看起来}相依。例如，我们可能在大量随机变量对中进行搜索，却没有进行多重检验校正。在这种情况下，关于变量存在相依关系的推断，并不能满足预期的第一类错误控制要求，见附录 A.2。
 \item 类似地，两个随机变量也可能都带有时间依赖，并遵循某种简单的物理规律，例如指数增长。此时，这些变量看起来彼此相依；但由于独立同分布假设被破坏，应用标准的独立性检验就缺乏依据。
\end{enumerate}
尤其是在报告随机变量之间的“伪相关”时，应牢记第 (2)、(3) 点；不少大众网站都会展示这类例子。
